\SourcePage{219}{222}
\section{Semiclassical motion in a centrally symmetric field}

In a centrally symmetric field, the particle's wave function separates, as
we know, into angular and radial parts. We first consider the former.

The dependence of the angular wave function on the angle $\varphi$ (fixed by
the quantum number $m$) is so simple that no question of finding approximate
formulas for it arises. Its dependence on the polar angle $\theta$, however,
is semiclassical by the general rule when the corresponding quantum number
$l$ is large (a more precise statement of this condition is given below).

We restrict the derivation of a semiclassical angular function to the case
most important in applications, namely states with zero magnetic quantum
number ($m=0$)\footnote{The opposite case, $m=l$, must in the limit correspond
to motion along a classical orbit lying in the equatorial plane
$\theta=\pi/2$. Indeed,
\[
 P_l^l(\cos\theta)=\operatorname{const}\cdot\sin^l\theta;
\]
as $l\to\infty$, this function (and hence also $|\psi|^2$) tends to zero at
every $\theta\ne\pi/2$.}. Apart from a constant factor, this function is the
Legendre polynomial $P_l(\cos\theta)$ (see (28.8)) and obeys
\begin{equation}
 \frac{d^2P_l}{d\theta^2}+\cot\theta\frac{dP_l}{d\theta}
 +l(l+1)P_l=0. \tag{49.1}
\end{equation}
The substitution
\begin{equation}
 P_l=\frac{\chi}{\sqrt{\sin\theta}} \tag{49.2}
\end{equation}
reduces it to the equation
\begin{equation}
 \chi''+\left[\left(l+\frac12\right)^2+
 \frac{1}{4\sin^2\theta}\right]\chi=0, \tag{49.3}
\end{equation}
which contains no first derivative and has the form of a one-dimensional
Schrödinger equation.

In (49.3), the role of the ``de Broglie wavelength'' is played by
$\lambda\sim1/l$.
\SourcePage{220}{223}
The requirement that $d\lambda/dx$ be small (condition (46.6)) gives
\begin{equation}
 \theta l\gg1,\qquad(\pi-\theta)l\gg1, \tag{49.4}
\end{equation}
which are the semiclassical conditions for the angular part of the wave
function. For large $l$ they hold through nearly the entire range of
$\theta$, failing only for angles very close to zero or to $\pi$.

When (49.4) holds, the second term in the square brackets in (49.3) may be
neglected against the first:
\[
 \chi''+\left(l+\frac12\right)^2\chi=0.
\]
The solution is
\begin{equation}
 \chi=\sqrt{\sin\theta}\,P_l(\cos\theta)
 =A\sin\left[\left(l+\frac12\right)\theta+\alpha\right] \tag{49.5}
\end{equation}
($A$ and $\alpha$ are constants).

For $\theta\ll1$, we may put $\cot\theta\approx1/\theta$ in (49.1); also
replacing $l(l+1)$ approximately by $(l+1/2)^2$, we obtain
\[
 \frac{d^2P_l}{d\theta^2}+\frac1\theta\frac{dP_l}{d\theta}
 +\left(l+\frac12\right)^2P_l=0,
\]
whose solution is a Bessel function of order zero:
\begin{equation}
 P_l(\cos\theta)=J_0\left[\left(l+\frac12\right)\theta\right],
 \qquad\theta\ll1. \tag{49.6}
\end{equation}
The constant factor has been set equal to unity because $P_l=1$ at
$\theta=0$. Approximation (49.6) is valid for all $\theta\ll1$. In
particular, it applies in the region $1/l\ll\theta\ll1$, where it must match
(49.5), valid for every $\theta\gg1/l$. When $\theta l\gg1$, the Bessel
function may be replaced by its large-argument asymptotic form, giving
\[
 P_l\approx\sqrt{\frac{2}{\pi l\theta}}
 \sin\left[\left(l+\frac12\right)\theta+\frac\pi4\right]
\]
(the $1/2$ may be neglected against $l$ in the coefficient). Comparison with
(49.5) gives $A=\sqrt{2/\pi l}$ and $\alpha=\pi/4$. We therefore obtain the
following final expression for $P_l(\cos\theta)$,
\SourcePage{221}{224}
applicable in the semiclassical case\footnote{Notice that it is precisely the
replacement of $l(l+1)$ by $(l+1/2)^2$ that has produced an expression
multiplied by $(-1)^l$ when $\theta$ is replaced by $\pi-\theta$, as required
for $P_l(\cos\theta)$.}:
\begin{equation}
 P_l(\cos\theta)=\sqrt{\frac{2}{\pi l\sin\theta}}
 \sin\left[\left(l+\frac12\right)\theta+\frac\pi4\right]. \tag{49.7}
\end{equation}
The normalized spherical function $Y_{l0}$ is consequently (cf. (28.8))
\begin{equation}
 Y_{l0}=\frac{1}{\pi\sqrt{\sin\theta}}
 \sin\left[\left(l+\frac12\right)\theta+\frac\pi4\right]. \tag{49.8}
\end{equation}

We now turn to the radial part of the wave function. It was shown in \secref{32}
that $\chi(r)=rR(r)$ obeys an equation identical to the one-dimensional
Schrödinger equation with potential energy
\[
 U_l(r)=U(r)+\frac{\hbar^2}{2mr^2}l(l+1).
\]
We may therefore use the results of the preceding sections, taking the
potential energy there to mean $U_l(r)$.

The case $l=0$ is simplest. There is no centrifugal energy, and if $U(r)$
satisfies the necessary condition (46.6), the radial wave function is
semiclassical throughout space. At $r=0$ we must have $\chi=0$, so the
semiclassical function $\chi(r)$ is determined in accordance with (47.6).

If $l\ne0$, the centrifugal energy must also satisfy (46.6). In the region of
small $r$ where the centrifugal energy has the order of the total energy,
$\lambda=\hbar/p\sim r/l$, and (46.6) gives $l\gg1$. Thus, when $l$ is not
large, the centrifugal energy violates the semiclassical condition in the
region of small $r$. It is readily verified that the correct phase of the
semiclassical wave function $\chi(r)$ is obtained by using the formulas for
one-dimensional motion with $l(l+1)$ in $U_l(r)$ replaced
\SourcePage{222}{225}
by $(l+1/2)^2$\footnote{Thus, in the simplest case of free motion ($U=0$),
the phase calculated from (48.1) with $U_l$ from (49.9) agrees at large $r$,
as it should, with the phase of (33.12).}:
\begin{equation}
 U_l(r)=U(r)+\frac{\hbar^2}{2mr^2}\left(l+\frac12\right)^2. \tag{49.9}
\end{equation}

The applicability of the semiclassical approximation to a Coulomb field
$U=\pm\alpha/r$ requires separate consideration. Of the entire region of
motion, the most significant part corresponds to distances at which
$|U|\sim|E|$, that is, $r\sim\alpha/|E|$. Semiclassical motion in this region
requires the reduced wavelength
$\lambdabar\sim\hbar/\sqrt{2m|E|}$ to be small compared with the size
$\alpha/|E|$ of the region. This gives
\begin{equation}
 |E|\ll\frac{m\alpha^2}{\hbar^2}, \tag{49.10}
\end{equation}
so the absolute value of the energy must be small compared with the particle's
energy in the first Bohr orbit. Condition (49.10) may also be written
\begin{equation}
 \frac{\hbar v}{\alpha}\gg1, \tag{49.11}
\end{equation}
where $v\sim\sqrt{|E|/m}$ is the particle's speed. Notice that this is the
reverse of condition (45.7) for applicability of perturbation theory to the
Coulomb field.

The region of small distances ($|U(r)|\gg E$) is of no interest in a
repulsive Coulomb field, since the semiclassical wave functions decay
exponentially when $U>E$. In an attractive field and for small $l$, however,
the particle can enter a region where $|U|\gg|E|$, raising the question of
the range of validity of the semiclassical approximation there. We use the
general condition (46.7), setting
\[
 F=-\frac{dU}{dr}=\frac{\alpha}{r^2},\qquad
 p\approx\sqrt{2m|U|}\sim\sqrt{\frac{m\alpha}{r}}.
\]
It follows that the range of validity is restricted to distances
\begin{equation}
 r\gg\frac{\hbar^2}{m\alpha}, \tag{49.12}
\end{equation}
that is, distances large compared with the ``radius'' of the first Bohr
orbit.
\SourcePage{223}{226}
\begin{center}\ProblemHeading\end{center}

Determine the behavior of the wave function near the origin if, as $r\to0$,
the field tends to infinity as $\pm\alpha/r^s$, with $s>2$.

\SolutionHeading For sufficiently small $r$, the wavelength is
\[
 \lambdabar\sim\frac{\hbar}{\sqrt{m|U|}}
 \sim\frac{\hbar r^{s/2}}{\sqrt{m\alpha}},
\]
and hence
\[
 \frac{d\lambdabar}{dr}\sim
 \frac{\hbar r^{s/2-1}}{\sqrt{m\alpha}}\ll1;
\]
thus the semiclassical condition is satisfied. In an attractive field,
$U\to-\infty$ as $r\to0$. The region near the origin is then classically
accessible, and the radial wave function $\chi\sim1/\sqrt p$, whence
\[
 \psi\sim r^{s/4-1}.
\]
In a repulsive field the region of small $r$ is classically inaccessible. The
wave function then tends exponentially to zero as $r\to0$. Omitting the
factor multiplying the exponential, we have
\[
 \psi\sim\exp\left[-\frac{2\sqrt{m\alpha}}{(s-2)\hbar}
 r^{-(s/2-1)}\right].
\]
